\vspace{5mm}
\begin{ledgroup}
\footnotesize
\pstart
\noindent
Bei den nachfolgend geschlossen edierten Stücken handelt es sich um
die in den \cite{01023}\textit{Acta Eruditorum} von März 1694 (S.\;112\,f.) 
anonym erschienene Rezension von 
\indexN{\textls{Vanni}, Giovanni Francesco 1638\textendash1709}\textsc{G.\,F. Vanni}, \cite{02001}\textit{Investigatio momentorum quibus gravia tendunt deorsum}, Rom 1693 (N.\;\ref{61037})
und um einen ursprünglich als Bestandteil derselben gedachten, zuletzt aber von Leibniz einbehaltenen (\glqq nicht abgangen\grqq)
Kommentar über \indexN{\textls{Vanni}, Giovanni Francesco 1638\textendash1709}Vannis Fehlschlüsse, der also nur noch als Abschrift im Konzept N.\;\ref{41824} erhalten ist.
\pend
\pstart
Während die gedruckte Rezension N.\;\ref{61037} neutral aus der \cite{02001}\textit{Investigatio} referiert,
setzt sich Leibniz in N.\;\ref{41824} mit \indexN{\textls{Vanni}, Giovanni Francesco 1638\textendash1709}Vannis Argumenten
auseinander und führt die eigene, bereits in früheren Texten (bes. N.\;\ref{49409} und N.\;\ref{41199}) geäußerte Kritik fort.
Eine Bemerkung am Anfang des Konzepts N.\;\ref{41824} gibt Auskunft über die gemeinsame Entstehung und das getrennte Schicksal der Stücke:
\glqq Haec sequentia addere constitueram, sed postea omisi, ne controversia relationi misceretur.\grqq\ 
Leibniz hatte den lediglich referierenden Teil der Rezension und die weiterführenden inhaltlichen Bemerkungen nicht nur gleichzeitig verfasst
\textendash\
in der Überschrift von N.\;\ref{41824} bezeichnet er die Rezension N.\;\ref{61037} als \glqq\textit{Actis} Lips. inserenda\grqq, also als fertiggestellt aber noch nicht erschienen \textendash\
sondern als zwei Teile ein und derselben Besprechung.
Obwohl Leibniz den Schutzes einer doppelten Anonymität genossen hätte
(denn in der rhetorischen Fiktion von N.\;\ref{41824} berichtet der \textit{anonyme} Rezensent von der kritischen Stellungnahme eines \textit{anonymen} Kritikers)
zog er es schließlich, im Interesse der Neutralität (\glqq ne controversia relationi misceretur\grqq) vor,
den Schlussteil des Textes auszulassen.
\pend
\pstart
Die Entstehung der Rezension lässt sich anhand des Briefwechsels mit 
\indexN{\textls{Mencke} (Menken, Menkenius, Menque), Otto 1644\textendash1707}O.~Mencke
datieren.
Mencke hatte Leibniz am 2.~(12.) Dezember 1693 ein Exemplar der \cite{02001}\textit{Investigatio} geschickt,
\glqq cum onere excerpendi\grqq\ und mit der Bitte, eine \glqq relation ad \textit{Acta}\grqq\ abzufassen 
(\cite{02097}\textit{LSB} I,~9 N.\;427, hier S.\;637).
Der genaue Zeitpunkt, zu dem Leibniz die fertige Rezension einsendete, ist nicht bekannt; 
in Frage kommen lediglich der 14.~(24.) Februar oder der 25.~Februar (7.~März) 1694, denn Mencke bedankt sich in seinem \cite{02098}Schreiben vom 7.~(17.) März für zwei entsprechend datierte, nach heutigem Kenntnisstand verschollene Leibnizbriefe
und für die Rezension, die schließlich im März 1694 (a.\:St.) in den \cite{01023}\textit{Acta} erscheint.
Die Abschrift (N.\;\ref{41824}) des Schlussteils der Besprechung entstand wohl gleichzeitig mit der Einsendung der gekürzten Rezension (N.\;\ref{61037}) und lässt sich daher etwa in der Zeitspanne zwischen den Briefen vom 24.~Februar und vom 7.~März datieren.
\pend
\pstart
Nach heutigem Wissenstand sind keine Notizen, Auszüge oder sonstige Spuren von Leibnizens Lektüre des Werkes überliefert; ebenfalls fehlen jegliche Konzepte, die die Entstehung der Rezension (und des einbehaltenen Kommentars) näher beleuchten könnten.
\pend 
\end{ledgroup}

